\begin{lem}\label{nsi} Let $G$ be a Polish group with a comeager conjugacy class. If $N$ is a normal subgroup of small index then $N=G$.
\end{lem}

\begin{proof}Let $\cC$ be the comeager conjugacy class. It suffices to show that $\cC\cap N\neq\emptyset$. Otherwise, by normality, $\cC\subset N$ and thus $G=\cC\cdot\cC\subset N$.

Since $N$ has small index then $N$ is not meager \cite[Th.\,4.1]{MR1231710} (see also \cite[Lem.\,6.8]{Kechris-Rosendal} for a very short proof), so $N\cap\cC\neq\emptyset$ and we are done.\end{proof}

Our proof that $G\wr\sinf$ has the small index subgroup property borrows the original ideas that lead to prove the property for $\sinf$ \cite[Th.\,1]{MR859950}. We not only use the result but also the proof itself and thus reproduce some of the arguments there. Let us recall that a moiety of $\Omega$ is a subset $\Sigma$ that is infinite and co-infinite.

\begin{proof}[Proof of Theorem~\ref{wr}] Let $H$ be a subgroup of $W$ of small index. Let $(\Sigma_i)$ be an infinite collection of disjoint moieties of $\Omega$. Let $W_i=G\wr_{\Sigma_i} \Sym(\Sigma_i)\leq\Sym(\Lambda\times \Omega)$ be the subgroup of permutations supported on $\Lambda\times\Sigma_i$. More precisely, $\left((g_\omega),\sigma\right)\in W_i$ if $\supp(\sigma)\subset \Sigma_i$ and $g_\omega=e$ for $\omega\notin \Sigma_i$. By disjointness of the supports, for $i\neq j$, $W_i\cap W_j$ is trivial and these two subgroups commute. Let $P$ be the subgroup $\Sym(\Lambda\times \Omega)$ isomorphic to $\prod_i W_i$. Let $H_i$ be the projection of $H\cap P$ on $W_i$. We have
\[
\prod_{i}|W_i\colon H_i|=\Bigl|P\colon \prod_{i}H_i\Bigr|\leq|P\colon H\cap P|\leq|W\colon H|<2^{\aleph_0}.
\]

This implies that for all but finitely many, $W_i=H_i$. We fix such an $i$ and simply note $\Sigma=\Sigma_i$ and $W'=W_i$. We denote by $G^\Sigma$ the subgroup of $G\wr\sinf$ with elements of the form $((g_\omega),e)$ and $g_\omega=e$ for $\omega\notin \Sigma$. Let $g\in H\cap G^\Sigma$ and $g'\in G^\Sigma$. There exists $h\in H\cap P$ such that $\pi_i(h)=g'$, where $\pi_i\colon P\to W_i$ is the projection. One has $g'g(g')^{-1}= hgh^{-1}\in H\cap G^\Sigma$. So, the subgroup $H\cap G^\Sigma$ is a normal subgroup of $G^\Sigma$ of small index. The group $G^\Sigma$ has a comeager conjugacy class because of \cite[Lem.\,11]{Rosendal-Solecki} and the fact that $G$ has a comeager conjugacy class. By Lemma~\ref{nsi}, $H\cap G^\Sigma=G^\Sigma$.

We denote by $\Sym(\Sigma)$ the subgroup of $G\wr\sinf$ of elements of the form $((g_\omega),\sigma)$, where $g_\omega=e$ for all $\omega\in\Omega$ and $\supp(\sigma)\subset \Sigma$. Since $\Sym(\Sigma)\cong\sinf$ and the non-trivial normal subgroups of $\sinf$ are the finitary symmetric and alternating subgroups, which are of index $2^{\aleph_0}$, we know that for $H\cap \Sym(\Sigma)=\Sym(\Sigma)$. Now, $W'=G\wr_\Sigma\Sym(\Sigma)$ is generated by $G^\Sigma$ and $\Sym(\Sigma)$. Thus $H\geq W'$.

Following the proof of \cite[Th.\,1]{MR859950}, we choose a continuum of almost disjoint moieties $(\Sigma_\alpha)$ of $\Sigma$ and an involution $g_\alpha\in\sinf$ exchanging $\Sigma_\alpha$ with $\Omega\setminus\Sigma$ and fixing pointwise $\Sigma\setminus\Sigma_\alpha$. By the pigeonhole principle, for some $\alpha\neq\beta$, $g_\alpha$ and $g_\beta$ are in the same $H$-class and thus $g=g_\beta^{-1} g_\alpha$ is in $H$. Let $\Sigma'=g(\Sigma)$, then $G\wr_{\Sigma'}\Sym(\Sigma')=gW'g^{-1}\leq H$. Since $\Sigma\cup\Sigma'=\Omega\setminus g_\beta(\Sigma_\alpha\cap\Sigma_\beta)$ and $\Sigma\cap\Sigma'=\Omega\setminus g_\beta(\Sigma_\alpha\cup\Sigma_\beta)$ are infinite, the first lemma of \cite{MR859950} shows that $H_0=G\wr_{\Sigma\cup\Sigma'}\Sym(\Sigma\cup\Sigma')\subset H$.

Let us denote by $F$ the finite set $g_\beta(\Sigma_\alpha\cap\Sigma_\beta)$. For $f\in F$, let $G_f\leq G\wr_\Omega\sinf$ be the corresponding copy of $G$ acting on $\Lambda\times\{f\}$. The subgroup $H_f=H\cap G_f$ has small index and thus is open in $G_f$. Now
\[
H\geq\left(\Pi_{f\in F}H_f\right)\times H_0.
\]
The right-hand side being open (containing the pointwise stabilizer of a finite number of points), $H$ is open as well.
\end{proof}
